\chapter{Spójność, przekroje i twierdzenie Mengera}
\label{ch:tg-menger}

Spójność grafu o niepustym zbiorze wierzchołków mówi,
czy każdą parę wierzchołków łączy choć jedna ścieżka. W zagadnieniach geometrycznych i kombinatorycznych
potrzebujemy często informacji silniejszej: \emph{ile wzajemnie
niezależnych ścieżek} łączy dwa punkty? Twierdzenie Mengera
wiąże odpowiedź z najmniejszą liczbą krawędzi albo wierzchołków,
które trzeba usunąć, aby przerwać połączenie.

\section{Przekroje krawędziowe i separatory wierzchołkowe}
\label{sec:tg-cuts}

Przekroje krawędziowe $\delta_G(S)$ rozumiemy zgodnie
z definicją \ref{def:tg-tree-edge-cut} w rozdziale
\ref{ch:tg-drzewa}. W niniejszym rozdziale, gdy mówimy
o rozdzieleniu grafu na dwie części, rozważamy nietrywialne
podziały, tj. $\varnothing\ne S\subsetneq V$.

\begin{definicja}[Separator krawędziowy]
Niech $G=(V,E)$ będzie skończonym grafem prostym.
Dla różnych wierzchołków $s,t\in V$ zbiór krawędzi $F\subseteq E$
nazywamy \emph{separatorem krawędziowym $s,t$}, jeśli
w grafie $G-F$ nie ma ścieżki z $s$ do $t$.
\end{definicja}

\begin{definicja}[Spójność krawędziowa i wierzchołkowa]
Niech $G=(V,E)$ będzie skończonym grafem prostym.
Jeżeli $G$ jest spójny i $|V|\ge2$, to
\emph{spójnością krawędziową} $\lambda(G)$ nazywamy
najmniejszą liczbę krawędzi, których usunięcie rozspójnia $G$.
\emph{Spójnością wierzchołkową} $\kappa(G)$ nazywamy
najmniejszą liczbę wierzchołków, których usunięcie
rozspójnia $G$ lub pozostawia graf z jednym wierzchołkiem.
Dla grafu niespójnego (w tym grafu zerowego)
przyjmujemy $\lambda(G)=\kappa(G)=0$,
a dla $K_1$ obie liczby definiujemy jako zero.
\end{definicja}

Warunek dotyczący pozostawienia jednego wierzchołka pozwala
otrzymać $\kappa(K_n)=n-1$ dla $n\ge1$.
Dla $n\ge2$ usunięcie mniej niż $n-1$ wierzchołków
pozostawia bowiem graf pełny z co najmniej dwoma
wierzchołkami, a taki graf jest spójny.

\begin{stwierdzenie}
\label{prop:tg-lambda-cut}
Niech $G=(V,E)$ będzie spójnym grafem prostym z $|V|\ge2$.
Wówczas
\[
\lambda(G)=\min_{\varnothing\ne S\subsetneq V}|\delta_G(S)|
\qquad\text{oraz}\qquad \lambda(G)\le \delta(G).
\]
\end{stwierdzenie}

\begin{proof}
Usunięcie wszystkich krawędzi $\delta_G(S)$ oddziela
$S$ od $V\setminus S$, więc
$\lambda(G)\le|\delta_G(S)|$.
Odwrotnie, jeśli $F$ rozspaja graf, wybieramy zbiór
wierzchołków jednej składowej $S$ grafu $G-F$.
Żadna krawędź $G$ między $S$ i jego dopełnieniem
nie mogła pozostać, więc $\delta_G(S)\subseteq F$.
Otrzymujemy nierówność przeciwną po minimalizacji.
Dla $S=\{v\}$ mamy $|\delta_G(S)|=d_G(v)$;
minimalizując stopień dostajemy $\lambda(G)\le\delta(G)$.
\end{proof}

\begin{definicja}[Separator wierzchołkowy]
Niech $s,t$ będą różnymi \emph{niesąsiadującymi} wierzchołkami
grafu prostego $G=(V,E)$. Zbiór $U\subseteq V\setminus\{s,t\}$
nazywamy \emph{separatorem wierzchołkowym $s,t$}, jeśli
w $G-U$ nie istnieje ścieżka z $s$ do $t$.
Liczbę $\kappa_G(s,t)$ definiujemy jako najmniejszą
moc takiego separatora. Minimum istnieje: zbiór
$V\setminus\{s,t\}$ jest separatorem, gdyż $st\notin E$.
Jeśli $s,t$ należą do różnych składowych, pusty zbiór
jest separatorem i $\kappa_G(s,t)=0$.
\end{definicja}

Warunek niesąsiadowania jest istotny: gdy $st\in E$,
nie istnieje żaden separator wierzchołkowy nieobejmujący
$s$ ani $t$, ponieważ bezpośrednia krawędź nadal łączy końce.
Dlatego wersję wierzchołkową twierdzenia Mengera
najpierw sformułujemy dla wierzchołków niesąsiadujących.

\begin{center}
\begin{tikzpicture}[scale=.92,line cap=round]
 \node[tgBlueVertex] (s) at (-3,0) {$s$};
 \node[tgTealVertex] (u) at (-1,.95) {$u$};
 \node[tgTealVertex] (v) at (-1,-.95) {$v$};
 \node[tgRedVertex] (x) at (1,.95) {$x$};
 \node[tgRedVertex] (y) at (1,-.95) {$y$};
 \node[tgBlueVertex] (t) at (3,0) {$t$};
 \draw[tgEdge] (s)--(u) (s)--(v) (u)--(x) (u)--(y)
       (v)--(x) (v)--(y) (x)--(t) (y)--(t);
 \draw[tgAccent,densely dashed,line width=1.6pt]
       (0,-1.48)--(0,1.48);
 \node[font=\footnotesize,text=tgCoral] at (0,1.76)
       {przekrój o czterech krawędziach};
\end{tikzpicture}

{\footnotesize Przekrój krawędziowy zależy od podziału wierzchołków,
a nie od sposobu przecięcia linii na rysunku.}
\end{center}

\section{Sieci przepływowe: pomocnicze narzędzie dowodowe}
\label{sec:tg-flows}

\begin{definicja}[Sieć i przepływ]
\emph{Siecią przepływową} nazywamy skończony graf skierowany
$D=(V,A)$ bez pętli, z dwoma różnymi wyróżnionymi wierzchołkami
$s,t$ i funkcją \emph{przepustowości}
$c:A\to\mathbb R_{\ge0}$.
\emph{Przepływem} jest funkcja $f:A\to\mathbb R_{\ge0}$
spełniająca $0\le f(a)\le c(a)$ dla wszystkich łuków oraz
\emph{równowagę} $\sum_{a=(u,v)}f(a)=\sum_{a=(v,w)}f(a)$
dla każdego $v\notin\{s,t\}$. Jego \emph{wartość} to
$|f|:=\sum_{a=(s,w)}f(a)-\sum_{a=(u,s)}f(a)$.
Dla $S\subseteq V$ z $s\in S$ i $t\notin S$ definiujemy
pojemność przekroju
$c(S,V\setminus S):=\sum_{(u,v)\in A,\,u\in S,\,v\notin S}c(u,v)$.
\end{definicja}

W tej i następnych definicjach sumujemy po łukach
wskazanej orientacji. Jeśli w sieci są oba łuki $u\to v$
i $v\to u$, traktujemy je jako różne elementy.
Symbol $|f|$ oznacza wartość przepływu, a nie jej wartość
bezwzględną. W szczególności dla dowolnego przepływu
dopuszczalnego wartość ta może być ujemna.

\begin{lemat}[Wartość przepływu przez przekrój]
\label{lem:tg-cutbound}
Niech $f$ będzie przepływem w skończonej sieci
$(D,s,t,c)$, gdzie $D=(V,A)$. Dla dowolnego $S\subseteq V$ z $s\in S$ i $t\notin S$
zachodzi
\[
|f|=\sum_{\substack{(u,v)\in A\\u\in S,\,v\notin S}}f(u,v)
 -\sum_{\substack{(u,v)\in A\\u\notin S,\,v\in S}}f(u,v)
 \le c(S,V\setminus S).
\]
\end{lemat}

\begin{proof}
Sumujemy równania zachowania przepływu po wszystkich
wierzchołkach $S\setminus\{s\}$ i dodajemy do nich
definicję wartości przepływu w $s$. Każdy łuk o obu końcach
w $S$ jest liczony raz z plusem i raz z minusem, więc się
redukuje. Pozostają jedynie łuki wychodzące z $S$
i wchodzące do $S$, z przeciwnymi znakami. Daje to równość.
Nierówność wynika z nieujemności przepływu na łukach
wchodzących i ograniczeń $f\le c$ na łukach wychodzących.
\end{proof}

\begin{definicja}[Sieć resztowa]
Niech $f$ będzie przepływem w sieci $(D,s,t,c)$.
Dla każdego oryginalnego łuku $a=(u,v)$
wprowadzamy dwie możliwości zmiany: przejście $u\to v$
o \emph{pojemności resztowej} $c(a)-f(a)$ oraz przejście
$v\to u$ o pojemności resztowej $f(a)$. Pomijamy
przejścia o zerowej pojemności. Różne łuki oryginalne
mogą dawać równoległe przejścia resztowe; pozostają one
rozróżnione. \emph{Ścieżką powiększającą} nazywamy
skierowaną ścieżkę z $s$ do $t$ w takiej sieci resztowej,
z parami różnymi wierzchołkami i z jawnym wyborem
identyfikatora każdego przejścia resztowego.
\end{definicja}

\begin{twierdzenie}[Maksymalny przepływ -- minimalny przekrój]
\label{thm:tg-maxflow}
Niech $(D,s,t,c)$ będzie skończoną siecią przepływową,
gdzie $D=(V,A)$, a $c:A\to\mathbb Z_{\ge0}$ jest funkcją
całkowitoliczbowych przepustowości. Wtedy istnieje
całkowitoliczbowy przepływ maksymalny i zachodzi
\[
\max_f |f|=\min_{\substack{s\in S\subseteq V\\t\notin S}}
 c(S,V\setminus S),
\]
gdzie maksimum przebiega po wszystkich dopuszczalnych
przepływach rzeczywistych w tej sieci.
\end{twierdzenie}

\begin{proof}[Dowód przez ścieżki powiększające]
Rozpoczynamy od przepływu zerowego. Dopóki istnieje
ścieżka powiększająca $P$ z $s$ do $t$, definiujemy
$\varepsilon$ jako najmniejszą dodatnią pojemność resztową
na tej ścieżce. Na oryginalnych łukach odpowiadających
krokom w przód zwiększamy przepływ o $\varepsilon$,
a na łukach odpowiadających krokom wstecz zmniejszamy
go o $\varepsilon$. Ponieważ ścieżka nie powtarza
wierzchołków, nie wykorzystuje obu przejść odpowiadających
temu samemu łukowi oryginalnemu ani tego samego przejścia
dwukrotnie. Każda zmiana jest zatem dopuszczalna na mocy
definicji $\varepsilon$: w przód mamy $f(a)+\varepsilon
\le c(a)$, a wstecz $f(a)-\varepsilon\ge0$.
Dla przejścia resztowego od $u$ do $v$ bilans
(odpływ minus dopływ) zwiększa się o $\varepsilon$
w $u$ i zmniejsza o $\varepsilon$ w $v$, niezależnie
od tego, czy zwiększamy przepływ na łuku $u\to v$,
czy zmniejszamy go na łuku $v\to u$.
W każdym wierzchołku wewnętrznym zmiany bilansu
się znoszą. Bilans w $s$ wzrasta o $\varepsilon$,
a bilans w $t$ maleje o $\varepsilon$. W szczególności
wartość przepływu zwiększa się dokładnie o $\varepsilon$.

Przy pojemnościach całkowitych również wszystkie przepływy
pośrednie i pojemności resztowe są całkowite, toteż
$\varepsilon\ge1$. Z lematu \ref{lem:tg-cutbound}
wartość przepływu jest ograniczona przez pojemność
dowolnego przekroju, a zatem algorytm zatrzyma się
po skończonej liczbie iteracji.

Niech $f$ będzie przepływem końcowym i niech $S$ będzie
zbiorem wierzchołków osiągalnych z $s$ po dodatnich
łukach resztowych. Skoro algorytm się zatrzymał, $t\notin S$.
Dla każdego oryginalnego łuku $u\to v$ z $u\in S$,
$v\notin S$, musi zachodzić $f(u,v)=c(u,v)$:
w przeciwnym razie istniałby dodatni resztowy krok
$u\to v$ i $v$ należałby do $S$.
Analogicznie dla oryginalnego łuku $v\to u$ z $v\notin S$,
$u\in S$ mamy $f(v,u)=0$, ponieważ jego dodatni przepływ
tworzyłby resztowy krok $u\to v$.
Z tożsamości przepływu przez przekrój otrzymujemy więc
$|f|=c(S,V\setminus S)$.

Lemat \ref{lem:tg-cutbound} zapewnia, że wartość żadnego
przepływu nie przekracza pojemności żadnego przekroju.
Znaleźliśmy przepływ i przekrój, dla których występuje
równość, zatem są one odpowiednio maksymalne i minimalne.
Każdy przepływ w konstrukcji jest całkowitoliczbowy.
\end{proof}

\begin{przyklad}[Rachunek przepływu i przekroju]
Rozważmy sieć z łukami $s\to a$, $s\to b$, $a\to t$,
$b\to t$ o pojemności $1$ oraz $a\to b$ o pojemności $2$.
Przepływ wysyłający jednostkę po drodze $s\to a\to t$
i jednostkę po drodze $s\to b\to t$ ma wartość $2$.
Przekrój $S=\{s\}$ ma pojemność $2$, więc na mocy
poprzedniego twierdzenia przepływ jest maksymalny.
Dodatkowy łuk $a\to b$ o większej pojemności
nie może zwiększyć przepływu: wąskim gardłem
są dwa łuki wychodzące z $s$.
\end{przyklad}

\begin{center}
\begin{tikzpicture}[scale=.95,>=Stealth,line cap=round]
 \node[tgBlueVertex] (s) at (-2.8,0) {$s$};
 \node[tgTealVertex] (a) at (0,1.15) {$a$};
 \node[tgTealVertex] (b) at (0,-1.15) {$b$};
 \node[tgBlueVertex] (t) at (2.8,0) {$t$};
 \draw[tgTree,-{Stealth[length=2.5mm]}] (s)--node[above,font=\footnotesize] {$1$}(a);
 \draw[tgTree,-{Stealth[length=2.5mm]}] (s)--node[below,font=\footnotesize] {$1$}(b);
 \draw[tgTree,-{Stealth[length=2.5mm]}] (a)--node[above,font=\footnotesize] {$1$}(t);
 \draw[tgTree,-{Stealth[length=2.5mm]}] (b)--node[below,font=\footnotesize] {$1$}(t);
 \draw[tgEdge,-{Stealth[length=2.5mm]}] (a)--node[right,font=\footnotesize] {$2$}(b);
 \draw[tgAccent,densely dashed] (-1.95,-1.60)--(-1.95,1.60);
 \node[text=tgCoral,font=\footnotesize] at (-1.95,1.95) {przekrój $S=\{s\}$};
\end{tikzpicture}

{\footnotesize Liczby przy łukach oznaczają przepustowości,
a zielone strzałki pokazują dwa niezależne jednostkowe przepływy.}
\end{center}

\section{Twierdzenie Mengera w wersji krawędziowej}
\label{sec:tg-menger-edges}

\begin{definicja}[Ścieżki krawędziowo rozłączne]
Niech $s,t$ będą różnymi wierzchołkami grafu prostego $G$.
Dwie ścieżki $s$--$t$ są \emph{krawędziowo rozłączne},
jeżeli nie mają żadnej wspólnej krawędzi; mogą natomiast
mieć wspólne wierzchołki pośrednie. Oznaczmy przez
$\rho_E(s,t)$ maksymalną liczność rodziny parami
krawędziowo rozłącznych ścieżek $s$--$t$.
Niech $\lambda_G(s,t)$ oznacza minimalną moc separatora
krawędziowego $s,t$.
\end{definicja}

\begin{twierdzenie}[Menger, wersja krawędziowa]
\label{thm:tg-menger-edge}
Dla dowolnego skończonego grafu prostego $G=(V,E)$
i różnych wierzchołków $s,t\in V$, przy oznaczeniach
z powyższej definicji, zachodzi
$\rho_E(s,t)=\lambda_G(s,t)$.
\end{twierdzenie}

\begin{proof}
Każdy separator krawędziowy przecina każdą ścieżkę
$s$--$t$. Jeżeli ścieżki są parami krawędziowo rozłączne,
to różne ścieżki potrzebują różnych krawędzi separatora.
Stąd $\rho_E(s,t)\le\lambda_G(s,t)$.

Dla nierówności odwrotnej zastępujemy każdą krawędź
$\{u,v\}$ dwoma przeciwnymi łukami $u\to v$, $v\to u$
o pojemności $1$. Każdy przekrój $S$ tej sieci ma
pojemność równą $|\delta_G(S)|$: dla każdej krawędzi
przekraczającej granicę $S$ dokładnie jeden z dwóch
łuków wychodzi z $S$. Zatem najmniejsza pojemność
przekroju jest równa $\lambda_G(s,t)$:
każdy taki przekrój jest separatorem, a dla dowolnego
separatora $F$ zbiór $S$ wierzchołków osiągalnych z $s$
w $G-F$ nie zawiera $t$ i spełnia $\delta_G(S)\subseteq F$.

Na mocy twierdzenia \ref{thm:tg-maxflow} istnieje
całkowitoliczbowy przepływ tej wartości. Jeśli na
obu przeciwnych łukach jakiejś krawędzi płynie
jednostka, zmniejszamy oba te przepływy do zera.
W każdym z dwóch końców krawędzi zmniejszają się
o jeden zarówno dopływ, jak i odpływ; bilans
w żadnym wierzchołku, w tym w $s$ i $t$, się nie
zmienia. Wartość przepływu pozostaje więc taka
sama, a ograniczenia przepustowości nadal zachodzą.
Wykonując tę operację dla wszystkich takich par
łuków, otrzymujemy przepływ, który wzdłuż każdej
krawędzi nieskierowanej jest jednokierunkowy albo
zerowy.

Rozłóżmy pozostały całkowity przepływ na ścieżki
i cykle, uzasadniając dokładnie ten krok. Jeśli w grafie
łuków o dodatnim przepływie istnieje skierowany cykl,
zmniejszamy przepływ na każdym jego łuku o najmniejszą
z występujących tam dodatnich wartości. Jest ona
całkowita, więc całkowitoliczbowość zostaje zachowana.
Każdy wierzchołek tego cyklu traci jednakową wielkość
dopływu i odpływu, zatem nie zmienia się żaden bilans,
ani wartość przepływu. Przepływ pozostaje dopuszczalny,
a co najmniej jeden dodatni łuk przestaje być dodatni.
Żaden nowy dodatni łuk nie powstaje. Po skończonej liczbie
takich operacjach pozostaje przepływ dodatni wyłącznie
na skierowanym grafie acyklicznym.
Jeśli jego wartość jest dodatnia, z $s$ wychodzi łuk
o dodatnim przepływie. Przechodzimy po takich łukach.
Zatrzymujemy się przy pierwszym osiągnięciu $t$.
Nie możemy wrócić do $s$ ani powtórzyć innego
wierzchołka, gdyż dodatnie łuki tworzą graf acykliczny.
Nie możemy też zakończyć w wierzchołku różnym od $s,t$:
dodatni dopływ i zachowanie przepływu wymuszają
tam dodatni odpływ. Skończoność grafu zapewnia więc
dotarcie do $t$. Na każdym łuku uzyskanej ścieżki
zmniejszamy przepływ o jeden. Całkowitoliczbowość
zapewnia nieujemność nowych wartości, równowaga
w wierzchołkach wewnętrznych pozostaje zachowana,
a wartość przepływu maleje dokładnie o jeden.
Dodatnie łuki nadal tworzą graf acykliczny, ponieważ
żaden nowy dodatni łuk nie powstaje. Powtarzając
ten krok, otrzymujemy dokładnie tyle ścieżek $s$--$t$,
ile wynosiła początkowa wartość $|f|$.
Pozostały przepływ o wartości zero nie jest potrzebny
do konstrukcji tych ścieżek.
Ponieważ każdy łuk ma przepływ co najwyżej $1$,
a przeciwne łuki tej samej krawędzi zostały skasowane,
otrzymane ścieżki nie mają wspólnych krawędzi.
Ich liczba jest równa wartości przepływu,
czyli $\lambda_G(s,t)$. To daje nierówność odwrotną.
\end{proof}

\begin{wniosek}
Niech $G$ będzie spójnym grafem prostym z co najmniej
dwoma wierzchołkami. Wtedy $\lambda(G)$ jest najmniejszą,
po wszystkich parach różnych wierzchołków $s,t$,
liczbą krawędziowo rozłącznych ścieżek łączących $s$ i $t$.
\end{wniosek}

\begin{proof}
Z twierdzenia \ref{thm:tg-menger-edge} liczba ścieżek
dla pary $s,t$ jest równa najmniejszej liczbie krawędzi,
których usunięcie rozdziela tę parę. Minimalizując po
wszystkich parach, otrzymujemy dokładnie najmniejszą
liczbę krawędzi rozspajających graf.
\end{proof}

\section{Twierdzenie Mengera w wersji wierzchołkowej}
\label{sec:tg-menger-vertices}

\begin{definicja}[Ścieżki wewnętrznie rozłączne]
Niech $s,t$ będą różnymi wierzchołkami grafu $G$.
Parami różne ścieżki $P_1,\dots,P_k$ z $s$ do $t$ nazywamy
\emph{wewnętrznie wierzchołkowo rozłącznymi}, jeżeli
żadne dwie z nich nie mają wspólnego wierzchołka
poza ewentualnie wspólnymi końcami $s,t$.
Niech $\rho_V(s,t)$ będzie największą możliwą
liczbą takich ścieżek. Dopuszczamy rodzinę pustą,
więc jeśli $s,t$ nie łączy żadna ścieżka, liczba ta
wynosi zero. Minimum separatorów wierzchołkowych
rozważamy tylko dla niesąsiadujących końców.
\end{definicja}

\begin{twierdzenie}[Menger, wersja wierzchołkowa]
\label{thm:tg-menger-vertex}
Niech $G=(V,E)$ będzie skończonym grafem prostym,
a $s,t\in V$ różnymi \emph{niesąsiadującymi} wierzchołkami.
Wtedy $\rho_V(s,t)=\kappa_G(s,t)$.
\end{twierdzenie}

\begin{proof}
Każdy separator wierzchołkowy $s,t$ musi zawierać
przynajmniej jeden wierzchołek wewnętrzny każdej
ścieżki $s$--$t$. Jeśli takie ścieżki są wewnętrznie
rozłączne, separator musi mieć co najmniej tyle
wierzchołków, ile jest ścieżek. Stąd
$\rho_V(s,t)\le\kappa_G(s,t)$.

Zbudujmy skierowaną sieć, w której \emph{rozdzielimy
wierzchołki na wejścia i wyjścia}. Dla każdego
$v\in V\setminus\{s,t\}$ tworzymy wierzchołki
$v^-$ i $v^+$ połączone łukiem $v^-\to v^+$
o pojemności $1$. Końce $s,t$ pozostawiamy
nierozdzielone. Każdą krawędź $uv\in E$
zastępujemy dwoma przeciwnymi łukami,
łącząc \emph{wyjście} jednego końca z \emph{wejściem}
drugiego: odpowiednio $u^+\to v^-$ i $v^+\to u^-$,
gdzie dla $s,t$ symbol $z^\pm$ oznacza po prostu $z$.
Każdemu z tych łuków nadajemy pojemność
$M:=|V|+1$.

Ponieważ $s$ i $t$ nie sąsiadują, usunięcie wszystkich
pozostałych wierzchołków oddziela $s$ od $t$.
W sieci wybierzmy zbiór
$S_0=\{s\}\cup\{v^-:v\in V\setminus\{s,t\}\}$.
Wszystkie łuki wychodzące z $s$ prowadzą do wejść
należących do $S_0$, a z każdego $v^-$ wychodzi
jedynie łuk $v^-\to v^+$.
Zatem łukami wychodzącymi z $S_0$ są dokładnie
$|V|-2$ łuki o pojemności $1$. Otrzymaliśmy
przekrój o pojemności $|V|-2<M$.
Każdy przekrój minimalny ma pojemność mniejszą niż $M$,
więc nie przecina żadnego z łuków o pojemności $M$.
W takim przekroju wszystkie łuki wychodzące
z jego zbioru $S$ są postaci $v^-\to v^+$
i mają pojemność $1$. Nie ma ograniczenia
na łuki wchodzące do $S$: nie wnoszą one
nic do pojemności przekroju.

Odpowiadające im wierzchołki tworzą separator
$s,t$ w grafie wyjściowym: gdyby po ich usunięciu
istniała ścieżka $s$--$t$, odpowiadająca jej droga
w sieci od $s$ do $t$ nie przecinałaby minimalnego
przekroju. Odwrotnie, dla dowolnego separatora
$U\subseteq V\setminus\{s,t\}$ możemy zbudować
przekrój sieci o pojemności co najwyżej $|U|$:
usuwamy łuki $v^-\to v^+$ dla $v\in U$ i bierzemy
zbiór wierzchołków osiągalnych z $s$ w pozostałej
sieci. W tej pozostałej sieci nie ma skierowanej
ścieżki od $s$ do $t$: przejście przez wewnętrzny
wierzchołek $v$ wymaga łuku $v^-\to v^+$,
więc taka ścieżka rzutowałaby się na spacer od $s$ do $t$
w $G-U$, wbrew własności separatora.
Zbiór osiągalny zawiera zatem $s$, lecz nie zawiera $t$.
Każdy łuk oryginalnej sieci wychodzący z tego zbioru
musiał zostać usunięty, gdyż inaczej jego koniec
byłby osiągalny. Są to łuki o pojemności $1$,
w liczbie co najwyżej $|U|$.
Stąd minimalna pojemność przekroju sieci jest
równa $\kappa_G(s,t)$.

Twierdzenie \ref{thm:tg-maxflow} daje całkowitoliczbowy
przepływ tej wartości. Rozkład przepływu na ścieżki
jednostkowe i cykle przeprowadzamy tak samo jak
w dowodzie wersji krawędziowej. Każda jednostkowa
ścieżka przechodząca przez wierzchołek $v\ne s,t$
musi użyć łuku $v^-\to v^+$, a jego pojemność wynosi $1$.
Zatem różne ścieżki nie mogą mieć wspólnego wierzchołka
wewnętrznego. Po zapomnieniu orientacji łuków
i utożsamieniu każdej pary $v^-,v^+$ z $v$
otrzymujemy ścieżki, a nie tylko spacery:
każde przejście przez $v$ zawiera oba wierzchołki
$v^-,v^+$ kolejno, a ścieżka w sieci nie powtarza
żadnego z nich. Końce $s,t$ także nie są powtarzane.
Otrzymujemy więc $\kappa_G(s,t)$ wewnętrznie rozłącznych
ścieżek $s$--$t$ w $G$. Wraz z pierwszą nierównością
to kończy dowód.
\end{proof}

\begin{center}
\begin{tikzpicture}[scale=.94,>=Stealth,line cap=round]
 \node[tgBlueVertex] (u) at (-2.4,0) {$u$};
 \node[tgTealVertex] (mi) at (-.25,0) {$v^-$};
 \node[tgRedVertex] (po) at (1.2,0) {$v^+$};
 \node[tgBlueVertex] (w) at (3.2,0) {$w$};
 \draw[tgEdge,-{Stealth[length=2.6mm]}] (u)--(mi);
 \draw[tgAccent,-{Stealth[length=2.6mm]}] (mi)--node[above=7pt,font=\footnotesize] {pojemność $1$}(po);
 \draw[tgEdge,-{Stealth[length=2.6mm]}] (po)--(w);
 \node[text=tgNavy,font=\footnotesize] at (.42,-.69)
    {przejście przez $v$ zużywa jedną jednostkę pojemności};
\end{tikzpicture}

{\footnotesize Rozcięcie wierzchołka na wejście i wyjście;
kolorowy łuk koduje ograniczenie użycia wierzchołka.}
\end{center}

\begin{przyklad}[Ścieżki rozłączne w prostym grafie]
Niech $G$ ma wierzchołki $s,t,a,b,c$ i krawędzie
$sa,at,sb,bt,sc,ct$. Istnieją trzy wewnętrznie
wierzchołkowo rozłączne ścieżki $s\! -\!a\! -\! t$,
$s\! -\! b\! -\! t$ i $s\! -\! c\! -\! t$.
Usunięcie $a,b,c$ rozdziela $s,t$, więc
$\rho_V(s,t)=\kappa_G(s,t)=3$.
Jeśli dodamy krawędź $ab$, liczba ta nie wzrośnie:
separator $\{a,b,c\}$ nadal ma moc trzy.
Przykład pokazuje, że dodatkowe krawędzie nie zawsze
zwiększają liczbę niezależnych tras.
\end{przyklad}

\begin{center}
\begin{tikzpicture}[scale=.95,line cap=round]
 \node[tgBlueVertex] (s) at (-2.6,0) {$s$};
 \node[tgBlueVertex] (t) at (2.6,0) {$t$};
 \foreach \n/\y/\col in {a/1.35/tgBlue,b/0/tgTeal,c/-1.35/tgCoral}{
   \node[tgVertex,draw=\col,fill=\col!12] (\n) at (0,\y) {$\n$};
   \draw[draw=\col,line width=2pt] (s)--(\n)--(t);
 }
 \node[font=\footnotesize,text=tgNavy] at (0,-1.98)
 {Trzy ścieżki rozłączne wewnętrznie: $s\to a\to t$, $s\to b\to t$, $s\to c\to t$.};
\end{tikzpicture}
\end{center}

\begin{wniosek}[Globalna wersja wierzchołkowa twierdzenia Mengera]
\label{cor:tg-menger-global-vertex}
Niech $G=(V,E)$ będzie skończonym spójnym grafem prostym
z $|V|\ge2$. Dla różnych $s,t\in V$ niech $\rho_V(s,t)$
oznacza największą liczność rodziny parami różnych,
wewnętrznie wierzchołkowo rozłącznych ścieżek od $s$ do $t$.
Wówczas
\[
\kappa(G)=\min_{\substack{s,t\in V\\s\ne t}}\rho_V(s,t).
\]
\end{wniosek}

\begin{proof}
Najpierw ustalmy dowolne różne $s,t\in V$.
Jeżeli $st\notin E$, to każdy separator wierzchołkowy
$s,t$ rozspaja cały graf, więc
$\kappa(G)\le\kappa_G(s,t)=\rho_V(s,t)$
na mocy twierdzenia~\ref{thm:tg-menger-vertex}.

Jeżeli $st\in E$, niech $H=G-st$ oznacza graf otrzymany
przez usunięcie tylko krawędzi $st$. Jedyną ścieżką
od $s$ do $t$ w $G$ wykorzystującą tę krawędź jest
ścieżka długości jeden: w przeciwnym razie któryś
z końców musiałby wystąpić ponownie. Każda inna
ścieżka należy do $H$. Ścieżkę długości jeden
możemy dołączyć do każdej rodziny wewnętrznie rozłącznych
ścieżek w $H$, gdyż nie ma wierzchołków wewnętrznych.
Stąd, stosując wersję lokalną w $H$, dostajemy
$\rho_V(s,t)=1+\kappa_H(s,t)$.
Wybierzmy separator $U$ dla $s,t$ w $H$
o mocy $\kappa_H(s,t)$. Jeśli $V\setminus U=\{s,t\}$,
to $\kappa(G)\le |V|-1=|U|+1$.
W przeciwnym razie wybierzmy
$w\in V\setminus(U\cup\{s,t\})$.
Jeśli $w$ należy do składowej $s$ w $H-U$, to
po usunięciu $U\cup\{s\}$ w $G$ wierzchołki $w,t$
leżą w różnych składowych: dodatkowa krawędź $st$
została usunięta wraz z $s$, a w $H-U$ nie istniała
ścieżka łącząca te składowe.
Jeśli $w$ nie należy do składowej $s$, to
po usunięciu $U\cup\{t\}$ wierzchołki $w,s$
leżą w różnych składowych z tego samego powodu.
W obu przypadkach $\kappa(G)\le |U|+1=\rho_V(s,t)$.
Udowodniliśmy dolne ograniczenie dla każdej pary.

Dla nierówności przeciwnej rozważmy najpierw graf
niepełny. Istnieją w nim niesąsiadujące wierzchołki,
a usunięcie pozostałych $|V|-2$ wierzchołków rozspaja
graf. Mamy więc $\kappa(G)\le |V|-2$.
Wybierzmy zbiór $U$ o mocy $\kappa(G)$ realizujący
minimum z definicji. Ponieważ po jego usunięciu
pozostają co najmniej dwa wierzchołki, graf $G-U$
musi być niespójny. Wybierzmy $s,t$ z różnych jego
składowych. Nie sąsiadują one w $G$, ponieważ ich
ewentualna krawędź pozostałaby w $G-U$.
Zbiór $U$ jest ich separatorem, więc
$\rho_V(s,t)=\kappa_G(s,t)\le |U|=\kappa(G)$.

Jeżeli $G=K_n$, gdzie $n=|V|\ge2$, to
$\kappa(G)=n-1$. Dla każdej pary $s,t$ bierzemy
ścieżkę $st$ oraz $n-2$ ścieżek $s,v,t$,
po jednej dla każdego $v\in V\setminus\{s,t\}$.
Są one parami wewnętrznie rozłączne. Więcej niż
$n-1$ takich ścieżek nie może istnieć, gdyż
poza jedyną ścieżką długości jeden każda z nich
zużywa co najmniej jeden z $n-2$ dostępnych
wierzchołków wewnętrznych. Zatem również w tym
przypadku minimum wynosi $\kappa(G)$.
\end{proof}

\section{Ile spójności zapewnia graf?}
\label{sec:tg-whitney}

\begin{twierdzenie}[Nierówności Whitneya]
\label{thm:tg-whitney}
Niech $G$ będzie skończonym spójnym grafem prostym
o co najmniej dwóch wierzchołkach. Wówczas
$\kappa(G)\le\lambda(G)\le\delta(G)$.
\end{twierdzenie}

\begin{proof}
Wybierzmy wierzchołek $v$ o stopniu minimalnym
$d_G(v)=\delta(G)$. Po usunięciu jego $\delta(G)$
sąsiadów $v$ jest izolowany: graf jest niespójny,
jeśli pozostał jeszcze inny wierzchołek, albo
jednoelementowy w przeciwnym razie. Z definicji
spójności wierzchołkowej w obu przypadkach
$\kappa(G)\le\delta(G)$.

Niech teraz $k=\lambda(G)$ i wybierzmy zbiór
$\varnothing\ne S\subsetneq V$ taki, że
$|\delta_G(S)|=k$; jego istnienie wynika ze
stwierdzenia \ref{prop:tg-lambda-cut}. Oznaczmy
$T=V\setminus S$ i $r=|S|$.

Jeśli $r>k$, usuńmy z $S$ wszystkie wierzchołki
będące końcami krawędzi przekroju $\delta_G(S)$.
Usunęliśmy co najwyżej $k$ wierzchołków,
w $S$ pozostał co najmniej jeden, a wszystkie
wierzchołki $T$ pozostały. Nie ma między nimi
żadnej krawędzi: każda taka krawędź wchodziłaby
do wybranego przekroju. Pozostały graf jest
niespójny, zatem $\kappa(G)\le k$.

Jeżeli $r\le k$, policzmy sumę stopni wierzchołków $S$.
Każdy z $r$ wierzchołków ma najwyżej $r-1$
sąsiadów w $S$, a łączna liczba krawędzi
z $S$ do $T$ wynosi $k$. Zatem
\[
r\delta(G)\le\sum_{v\in S}d_G(v)\le r(r-1)+k,
\qquad
\delta(G)\le r-1+\frac{k}{r}\le k.
\]
Ostatnia nierówność wynika z $1\le r\le k$:
po pomnożeniu przez $r$ jest równoważna
$(r-1)(r-k)\le0$. Z pierwszej części dowodu
$\kappa(G)\le\delta(G)\le k=\lambda(G)$.
Pozostała nierówność $\lambda(G)\le\delta(G)$
wynika ze stwierdzenia \ref{prop:tg-lambda-cut}.
\end{proof}



\begin{przyklad}
Dla cyklu $C_n$ przy $n\ge3$ mamy
$\delta(C_n)=2$. Usunięcie jednej krawędzi pozostawia
ścieżkę rozpinającą, a usunięcie jednego wierzchołka
pozostawia spójną ścieżkę z co najmniej dwoma
wierzchołkami. Stąd $\lambda(C_n),\kappa(C_n)\ge2$,
a z nierówności Whitneya obie liczby są równe $2$.

Dla drzewa $G$ z co najmniej dwoma wierzchołkami
każda krawędź jest mostem, więc $\lambda(G)=1$.
Drzewo jest spójne i ma co najmniej dwa wierzchołki,
zatem $\kappa(G)\ge1$; nierówności Whitneya dają
$\kappa(G)=1$. Ponieważ drzewo ma liść, także
$\delta(G)=1$.

Dla grafu pełnego $K_n$ z $n\ge2$ już wykazaliśmy
$\kappa(K_n)=n-1$, a $\delta(K_n)=n-1$.
Nierówności Whitneya dają więc również
$\lambda(K_n)=n-1$.
\end{przyklad}

\begin{przyklad}[Ostre nierówności Whitneya]
\label{ex:tg-whitney-strict}
Niech $G$ powstanie z dwóch rozłącznych kopii $K_4$
o zbiorach wierzchołków $A=\{a_1,a_2,a_3,a_4\}$
i $B=\{b_1,b_2,b_3,b_4\}$ przez dodanie wierzchołka
$z$ oraz czterech krawędzi $za_1,za_2,zb_1,zb_2$.
Nie dodajemy innych krawędzi. Wykażemy, że
$\kappa(G)=1<\lambda(G)=2<\delta(G)=3$.
Graf jest spójny, lecz usunięcie $z$
pozostawia dwie rozłączne kopie $K_4$,
więc $\kappa(G)=1$.
Usunięcie $za_1,za_2$ oddziela $A$ od reszty grafu,
zatem $\lambda(G)\le2$. Żadna pojedyncza krawędź
nie jest jednak mostem. Każda krawędź wewnątrz
$A$ lub $B$ leży na trójkącie w odpowiedniej
kopii $K_4$. Krawędzie $za_1,za_2$ leżą na
trójkącie o wierzchołkach $z,a_1,a_2$,
a krawędzie $zb_1,zb_2$ na trójkącie
o wierzchołkach $z,b_1,b_2$. Usunięcie krawędzi
leżącej na cyklu nie rozspaja grafu, więc
$\lambda(G)\ge2$.
Wreszcie wierzchołki $a_3,a_4,b_3,b_4$ mają stopień
$3$, a pozostałe wierzchołki, w tym $z$, mają stopień
$4$. Otrzymujemy $\delta(G)=3$.

Dla dwóch trójkątów o wspólnym wierzchołku mamy $\kappa=1$
(usunięcie tego wierzchołka) oraz $\lambda=\delta=2$
(brak mostów i stopień minimalny $2$).
Dwie kopie $K_4$ połączone jednym mostem mają
$\kappa=\lambda=1$ (spójność i nierówności Whitneya)
oraz $\delta=3$.
\end{przyklad}
